\section[Lie algebra $W(2,2)$ and the twisted Heisenberg--Virasoro Lie algebra at level zero]{Lie algebra $\boldsymbol{W(2,2)}$ and the twisted Heisenberg--Virasoro\\ Lie algebra at level zero}\label{rep-th-review}
$W(2,2)$ is a Lie algebra with basis $\{L(n), W(n),C_{L},C_{W}\colon n\in\mathbb{Z\}}$ over $\C$, and a Lie bracket
\begin{gather*}
[ L(n) ,L(m) ] =(n-m) L(n+m) +\delta_{n,-m}\frac{n^{3}-n}{12}C_{L}, \\ %\label{1} \\
[ L(n) ,W(m) ] =(n-m)W(n+m) +\delta_{n,-m}\frac{n^{3}-n}{12}C_{W}, \\
[ W(n) ,W(m) ] =[ \cdot,C_{L}] =[ \cdot,C_{W}] =0.
\end{gather*}

Highest weight representation theory over $W(2,2)$ was studied in~\cite{JP,R-JMP}. However, representations treated in these papers have
equal central charges $C_{L}=C_{W}$. These results have recently been generalised to $C_{L}\neq C_{W}$ in~\cite{JZ}. Here we state the most important results. Verma module with central charge $(c_{L},c_{W}) $ and highest weight $(h,h_{W}) $ is denoted by $V^{W(2,2)}( c_{L},c_{W},h,h_{W})$, its highest weight vector
by $v_{h,h_W}$ and irreducible quotient module by $L^{W(2,2)} (c_{L},c_{W},h,h_{W})$.

Recall the def\/inition of a cosingular vector. Homogeneous vector $v\in M$ is called cosingular (or subsingular) if it is not singular in $M$ and if there
is a proper submodule $N \subset M$ such that $v+N$ is a singular vector in~$M/N$.

\begin{Theorem}[{\cite{JZ, R-JMP}}]\label{W-struktura}Let $c_{W}\neq0$.
\begin{enumerate}\itemsep=0pt
\item[$(i)$]Verma module $V^{W(2,2)}(c_{L},c_{W},h,h_{W}) $ is reducible if and only if $h_{W}=\frac{1-p^{2}}{24}c_{W}$ for some $p\in \N$. In that case, there exists a singular vector $u_p^{\prime}\in \C [ W(-1) ,\ldots,W(-p)]v_{h,h_W} $ such that $U(W(2,2)) u_p^{\prime} \cong V^{W(2,2)}(c_{L},c_{W},h+p,h_{W}) $.

\item[$(ii)$]A quotient module\footnote{This module is denoted by $L^{\prime}$ in \cite{JZ,R-JMP}. We
change notation to $\widetilde{L}$ due to use of superscript $W(2,2)$.}
\begin{gather*}V^{W(2,2)}(c_{L},c_{W},h,h_{W}) / U(W(2,2)) u_p^{\prime} =: \widetilde{L}^{W(2,2)}( c_{L},c_{W},h_{p,r},h_{W})
\end{gather*}
is reducible if and only if $ h = h_{p,r} $ for some $r\in\N$. In that case, there is a cosingular vector $%
u_{rp}\in V^{W(2,2)}(c_{L},c_{W},h,h_{W}) _{h+rp}$ such that $\overline{u_{rp}} := u_{rp} + U(W(2,2)) u_p^{\prime}$ is a~singular vector in $\widetilde{L}^{W(2,2)}( c_{L},c_{W},h_{p,r},h_{W})$ which generates a submodule isomorphic to $L^{W(2,2)}( c_{L},c_{W},h_{p,r}+rp,h_{W})$. The short sequence
\begin{align}
0 &\rightarrow L^{W(2,2)}(c_{L},c_{W},h_{p,r}+rp,h_{W}) \rightarrow\widetilde{L}^{W(2,2)} ( c_{L},c_{W},h_{p,r},h_{W} ) \nonumber \\
&\rightarrow L^{W(2,2)} ( c_{L},c_{W},h_{p,r},h_{W} ) \rightarrow 0, \label{seq}
\end{align}
where the highest weight vector in $L^{W(2,2)}(c_{L},c_{W},h_{p,r}+rp,h_{W})$ maps to $\overline{u_{rp}}$ is exact.
\end{enumerate}
\end{Theorem}

Def\/ine
\begin{gather*}
{\mathcal {AT} }_{W(2,2) } ( c_L, c_{W} ) = \left\{ \left(h_{p,r}, \frac{1-p^2}{24}c_W\right) |\, p,r \in \N \right\}.
\end{gather*}
\begin{Remark}
We will refer to the (modules of) highest weights $(h, h_W) \in {\mathcal {AT} }_{W(2,2) } ( c_L, c_{W} ) $ as \textit{atypical} for $W(2,2)$, and otherwise as \textit{typical}. Again, we refer to a highest weight $W(2,2)$-module as (a)typical depending on its highest weight. So a Verma module over $W(2,2)$ contains a nontrivial cosingular vector if and only if it is atypical.
\end{Remark}

\begin{Proposition} \label{svi cosing}
Let $h_W=\frac{1-p^2}{24}c_W$, $p\in\N$.
\begin{enumerate}\itemsep=0pt
\item[$(i)$]Let $(h_{p,r},h_W)$, $r\in\N$ be an atypical weight and $k\in\Z$. Then $(h_{p,r}+kp,h_W)$ is atypical if and only if $k<\frac{r}{2}$.
\item[$(ii)$]Atypical Verma module $V^{W(2,2)}(h_{p,r},h_W)$ contains exactly $\lfloor \frac{r+1}{2} \rfloor$ cosingular vectors. The weights of these vectors are $h_{p,r}+(r-i)p = h_{p,-r+2i}$, $i=0,\ldots,\lfloor \frac{r-1}{2} \rfloor$.
\end{enumerate}
\end{Proposition}

\begin{proof}
(i) Directly from Theorem~\ref{W-struktura} since $h_{p,r}+kp = h_{p,r-2k}$.

(ii) Follows from~(i) since $V^{W(2,2)}(h_{p,r},h_W)$ contains an inf\/inite chain of submodules isomorphic to Verma modules of highest weights $h_{p,r}+ip = h_{p,r-2i}$, $i>0$. Applying Theorem~\ref{W-struktura} to each of these submodules we obtain cosingular vectors of weights
\begin{gather*}
h_{p,r-2i}+(r-2i)p = h_{p,r}+(r-i)p = h_{p,-r+2i}
\end{gather*}
as long as $r-2i>0$.
\end{proof}

\begin{Remark}\label{pbw}
Standard PBW basis for $V^{W(2,2)}(c_{L},c_{W},h,h_{W}) $ consists of vectors
\begin{gather*}
W(-m_{s}) \cdots W(-m_{1}) L(-n_{t}) \cdots L ( -n_{1} ) v_{h,h_W}
\end{gather*}
such that $ m_{s}\geq\cdots\geq m_{1}\geq1$, $n_{t}\geq\cdots\geq n_{1}\geq1 $. The only nonzero component of $u_{rp}$ belonging to $\C [ L(-1) ,L ( -2), \ldots ] v$ is $L(-p)^{r}v_{h,h_W}$~\cite{R-JMP}.
\end{Remark}

Def\/ine $P_{2}(n) =\sum\limits_{i=0}^{n}P(n-i)P(i)$ where $P $ is a partition function with $P(0)=1$. We have the following character
formulas~\cite{R-JMP}
\begin{gather*}
\ch V^{W(2,2)}( c_{L},c_{W},h,h_{W}) =q^{h}\sum_{n\geq0}P_{2}(n)q^{n}=q^{h}\prod\limits_{k\geq1}\big(1-q^{k}\big)^{-2},
\end{gather*}
for all $h,h_W \in \C$. If $h_W = \frac{1-p^2}{24}c_W$, then
\begin{gather*}
\ch \widetilde{L}^{W(2,2)}(c_{L},c_{W},h,h_{W}) =q^{h}\big(1-q^{p}\big)\sum_{n\geq0}P_{2}(n)q^{n}
 =q^{h}\big(1-q^{p}\big)\prod\limits_{k\geq1}\big(1-q^{k}\big)^{-2}.
\end{gather*}
If $(h,h_W)$ is typical for $W(2,2)$, then this is the character of an irreducible highest weight module. Finally, the character of atypical irreducible module is
\begin{align*}
\ch L^{W(2,2)}( c_{L},c_{W},h_{p,r},h_{W})& =q^{h_{p,r}}\big(1-q^{p}\big)\big(1-q^{rp}\big)\sum_{n\geq0}P_{2}(n)q^{n} \\
& =q^{h_{p,r}}\big(1-q^{p}\big)(1-q^{rp})\prod\limits_{k\geq1}\big(1-q^{k}\big)^{-2}.
\end{align*}

The twisted Heisenberg--Virasoro algebra $\HVir$ is the universal central extension of the Lie algebra of dif\/ferential operators on a circle of order at most one. It is the inf\/inite-dimensional complex Lie algebra with a basis
\begin{gather*}
\{L(n),I(n)\colon n\in\Z\}\cup\{C_{L},C_{LI},C_{I}\}
\end{gather*}
and commutation relations
\begin{gather*}
[ L(n),L(m)] =(n-m)L(n+m)+\delta_{n,-m}\frac{n^{3}-n}{12}C_{L},\\
[ L(n),I(m)] =-mI(n+m)-\delta_{n,-m}\big(n^{2}+n\big)C_{LI}, \\
[ I(n),I(m)] =n\delta_{n,-m}C_{I}, \qquad
[ {\HVir},C_{L}] =[ {\HVir},C_{LI}] =[ \HVir,C_{I}] =0.
\end{gather*}
The Lie algebra $\HVir$ admits the following triangular decomposition
\begin{gather}
\HVir =\HVir ^{-} \oplus \HVir^0 \oplus \HVir^{+}, \label{triangular} \\
\HVir ^{\pm} = \operatorname{span}_{\mathbb C} \{ I( \pm n), L(\pm n) \,|\, n \in {\Z}_{>0}\}, \qquad
\HVir^0 = \operatorname{span}_{\mathbb C} \{ I(0), L(0), C_L, C_{L,I}, C_{I} \}. \nonumber
\end{gather}

Although they seem to be two similar extensions of the Virasoro algebra, representation theories of $W(2,2)$ and $\HVir$ are dif\/ferent. The main reason for that lies in the fact that $I(0) $ is a~central element, while $W(0) $ is not. However, applying free f\/ield realization, we shall see that highest weight modules over the two algebras are related.

Denote by $V^{\HVir} ( c_{L},c_{I},c_{L,I},h,h_{I} ) $ the Verma module and by $v_{h,h_I}$ its highest weight vector. $C_{L},$ $C_{I},$ $C_{L,I},$ $L(0) $ and $I(0) $ act on $v_{h,h_I}$ by scalars $c_{L}$, $c_{I}$, $c_{L,I}$, $h$ and $h_{I}$, respectively. Then $\left(c_{L,}c_{I},c_{L,I}\right) $ is called a central charge, and $(h,h_{I}) $ a highest weight. In this paper we consider central charges $(c_{L},0,c_{L,I}) $ such that $c_{L,I}\neq0$.

\begin{Theorem}[\cite{Billig}]\label{billig}
Let $c_{L,I}\neq0$. Verma module $V^{\HVir}(c_{L},0,c_{L,I},h,h_{I}) $ is reducible if and only if $h_{I}=(1\pm p) c_{L,I}$ for some $p\in\N$. In that case, there is a singular vector $v_p^{\pm}$ of weight $p$, which generates a maximal submodule in $V^{\HVir} ( c_{L},0,c_{L,I},h,h_{I} ) $ isomorphic to $V^{\HVir} ( c_{L},0,c_{L,I},h+p,h_{I} ) $.
\end{Theorem}

\begin{Remark} \label{kompon}
In case $h_{I}=(1+p) c_{L,I}$ an explicit formula for a singular vector $v_p^+$ is obtained using Schur polynomials in $I(-1), \ldots, I(-p)$. See~\cite{nas} for details. Assume that $x \in U(W(2,2))_-$ is such that $xv_{h,h_I} \in V^{\HVir} ( c_{L},0,c_{L,I},h,h_{I} ) $ lies in a~maximal submodule. Then $x$ does not have a nontrivial additive component (in PBW basis) that belongs to $\C [ L(-1) ,L(-2), \ldots ]$~\cite{Billig}.
\end{Remark}

There is an inf\/inite chain of Verma submodules generated by singular vectors $v_{kp}^{\pm}$, $k\in\N$, with all the subquotients being irreducible. Note that there is no mention of $\widetilde{L}^{\HVir}$ since there are no cosingular vectors in $V^{\HVir}$.

The following character formulas were obtained in \cite{Billig}:
\begin{gather*}
\ch V^{\HVir}( c_{L},0,c_{L},h,h_{I}) =q^{h}\sum_{n\geq0}P_{2}(n)q^{n}=q^{h}\prod\limits_{k\geq1}\big(1-q^{k}\big)^{-2}, \\
\ch L^{\HVir} ( c_{L},0,c_{L},h,h_{I} ) =q^{h}\big(1-q^{p}\big)\sum_{n\geq0}P_{2}(n)q^{n} =q^{h}\big(1-q^{p}\big)\prod\limits_{k\geq1}\big(1-q^{k}\big)^{-2}.
\end{gather*}

\begin{Remark}
Throughout the rest of the paper we work with highest weight modules over the Lie algebras $W(2,2)$ and $\HVir$ so we always denote algebra in superscript. In order to avoid too cumbersome notation, we omit central charges. Therefore, we write $V^{\HVir}(h,h_{I}) $ for Verma module over~$\HVir$, ${V}^{W(2,2)}(h,h_{W}) $ for Verma module over $W(2,2)$ and so on. We always assume that $c_{W}$ and $c_{L,I}$ are nonzero. Moreover, if we work with several modules over both algebras, $c_{L}$ is equal for all modules.

We shall write $\langle x \rangle_{W(2,2)}$ for a cyclic submodule $U (W(2,2))x$ and $\langle x \rangle_{\HVir}$ for $U(\HVir)x$. Finally, $\cong_{W(2,2)}$ denotes an isomorphism of $W(2,2)$-modules.
\end{Remark}

\section{Irreducible highest weight modules}\label{result}
In this section we present main results of the paper which completely describe the structure of (irreducible) highest weight modules for $\HVir$ as $W(2,2)$-modules. The main tool is the homomorphism between $W(2,2)$ and the Heisenberg--Virasoro vertex algebras from~\cite{nas}.

$L^{W(2,2) }( c_{L},c_{W},0,0) $ is a simple universal vertex algebra associated to Lie algebra $W(2,2)$ (cf.\ \cite{R-JMP, Zhang-Dong}) which we denote by $L^{W(2,2)}(c_{L},c_{W})$. It is generated by f\/ields
\begin{gather*}
L(z)=Y(\omega ,z)=\sum_{n\in {\Z}}L(n)z^{-n-2},\qquad W(z)=Y(W,z)=\sum_{n\in {\Z}}W(n)z^{-n-2},
\end{gather*}
where $\omega =L(-2)\mathbf{1}$ and $W=W(-2)\mathbf{1}$. Each highest weight $W(2,2) $-module is also a module over a vertex operator algebra $L^{W(2,2) }(c_{L},c_{W})$.

Likewise (see \cite{Billig3}) $L^{\HVir}(c_{L},0,c_{L,I},0,0) $ is a simple Heisenberg--Virasoro vertex operator algebra, which we denote by $L^{\HVir}(c_{L},c_{L,I}) $. This algebra is generated by the f\/ields
\begin{gather*}
L(z)=Y(\omega ,z)=\sum_{n\in {\Z}}L(n)z^{-n-2},\qquad I(z)=Y(I,z)=\sum_{n\in {\Z}}I(n)z^{-n-1},
\end{gather*}
where $\omega =L(-2)\mathbf{1}$ and $I=I(-1)\mathbf{1}$. Moreover, highest weight $\HVir$-modules are modules over a~vertex operator algebra
$L^{\HVir}(c_{L,}c_{L,I})$.

It was shown in \cite{nas} that there is a monomorphism of vertex operator algebras
\begin{align}
\Psi \colon \ L^{W(2,2)}(c_{L},c_{W})& \rightarrow L^{\HVir}(c_{L},c_{L,I}),\label{Psi} \\
\omega & \mapsto L(-2)\mathbf{1}, \notag \\
W& \mapsto (I(-1)^{2}+2c_{L,I}I(-2))\mathbf{1}, \notag
\end{align}%
where $c_{W}=-24c_{L,I}^{2}$. By means of $\Psi $, each highest weight module over $\HVir$ becomes an \linebreak $L^{W(2,2)}(c_{L},c_{W})$-module and
therefore a module over $W(2,2) $. In particular, $\Psi $ induces a non-trivial $W(2,2)$-homomorphism (which we shall denote by the same letter)
\begin{gather*}
\Psi \colon \ V^{W(2,2)}(c_{L},c_{W},h,h_{W}) \rightarrow V^{\HVir} ( c_{L},0,c_{L,I},h,h_{I} ),
\end{gather*}
where $c_{W}=-24c_{L,I}^{2}$ and $h_{W}=h_{I}(h_{I}-2c_{L,I})$. $\Psi$ maps the highest weight vector $v_{h,h_{W}}$ to the highest weight vector
$v_{h,h_{I}}$ and the action of $W(-n)$ on $V^{\HVir} ( c_{L},0,c_{L,I},h,h_{I} )$ is given by
\begin{gather}
W(-n) \equiv 2c_{L,I}(n-1) I(-n)+\sum_{i\in \Z}I(-i)I(-n+i),\label{W u I} \\
W(-n) \equiv 2 c_{L,I}\left(n- 1 + \frac{h_I}{c_{L,I} } \right) I(- n)+\sum_{i\ne 0,n } I(-i)I(- n+i).\nonumber
\end{gather}

Note that $h_{W}=\frac{1-p^{2}}{24}c_{W}$ if and only if $h_{I}= ( 1\pm p) c_{L,I}$, so either both of these Verma modules are irreducible, or they are reducible with singular vectors at equal levels. Moreover, $ (h,h_W) \in {\mathcal {AT} }_{W(2,2)} ( c_L, c_W ) $ if and only if $(h,h_I) \in {\mathcal {AT} }_{\HVir} ( c_L, c_{L,I} )$.

Throughout the rest of this section we assume that $ c_W = -24c_{L,I}^2 $.

\begin{Lemma}[{\cite[Lemma 7.2]{nas}}] \label{psi}Suppose that $h_{I}\neq(1-p) c_{L,I}$ for all $p\in \N$. Then $\Psi$ is an isomorphism of $W(2,2)$-modules. In particular, if $ h_I \neq (1\pm p) c_{L,I} $ for $p \in \N $, then
\begin{gather*}
L^{\HVir}(h,h_I)\cong_{W(2,2)}L^{W(2,2)}(h,h_W),
\end{gather*}
where $h_W = h_I (h_I-2c_{L,I}) $.
\end{Lemma}

\begin{Lemma}\label{H-sing je W-sing}Suppose that $x\in V^{\HVir}(h,h_I)$ is $\HVir$-singular. Then $x$ is $W(2,2)$-singular as well. In particular, if $y$ is a homogeneous vector such that $x = \Psi(y)$, then $y$ is either singular or cosingular vector in $V^{W(2,2)}(h,h_W)$.
\end{Lemma}

\begin{proof}
Let $x\in V^{\HVir}(h,h_I)$ be a $\HVir$-singular vector, i.e., $L(k)x=I(k)x=0$ for all $k\in\N$. From (\ref{W u I}) we have
\begin{gather*}
W(n)x = -2c_{L,I}(n+1) I(n)x+\sum_{i\in \Z}I(-i)I(n+i)x,
\end{gather*}
so $W(n)x=0$ for all $n\in\N$. Therefore, $x$ is $W(2,2)$-singular.
If $x=\Psi(y)$, then $L(k)y, W(k)y$ $\in \Ker \Psi$ for $k>0$. Therefore $y + \Ker \Psi$ is a singular vector in $V^{W(2,2)}(h,h_W) / \Ker \Psi$.
\end{proof}

\begin{Theorem} \label{t33} \label{L za +}Let $p\in\N$.
\begin{enumerate}\itemsep=0pt
\item[$(i)$] If $(h,(1+p) c_{L,I})$ is typical for $\HVir$ $($equivalently if $\big(h,\frac{1-p^{2}}{24}c_{W}\big) $ is typical for $W(2,2))$ then
\begin{gather}
L^{\HVir}(h,(1+p) c_{L,I}) \cong_{W(2,2) }L^{W(2,2)}\left( h,\frac{1-p^{2}}{24}c_{W}\right).\label{L + gen}
\end{gather}

\item[$(ii)$] If $(h_{p,r},(1+p) c_{L,I}) \in {\mathcal {AT} }_{\HVir} ( c_L, c_{L,I} )$ $($equivalently if $ (h_{p,r},\frac{1-p^{2}}{24}c_{W}) \in {\mathcal {AT} }_{W(2,2)} ( c_L, c_W ))$ then
\begin{gather*}
L^{\HVir}(h_{p,r},(1+p) c_{L,I}) \cong_{W(2,2) }\widetilde{L}^{W(2,2)}\left( h_{p,r},\frac{1-p^{2}}{24}c_{W}\right) %\label{L + spec}
\end{gather*}
and the short sequence of $W(2,2)$-modules
\begin{align}
0 & \rightarrow L^{W(2,2)}\left( h_{p,r}+rp,\frac{1-p^{2}}{24}c_{W}\right)\rightarrow L^{\HVir}(h_{p,r},(1+p) c_{L,I}) \label{seq +} \\
& \rightarrow L^{W(2,2)}\left( h_{p,r},\frac{1-p^{2}}{24}c_{W}\right) \rightarrow0 \nonumber
\end{align}
is exact.
\end{enumerate}
\end{Theorem}

\begin{proof}
By Lemma \ref{psi}, $\Psi$ is an isomorphism of Verma modules and thus by Lemma \ref{H-sing je W-sing} it maps a $W(2,2) $-singular vector $u^{\prime}_p$ to an $\HVir$-singular vector $v_p^+$. If $h\neq h_{p,r}$, both of these vectors generate maximal submodules in respective Verma modules so~(\ref{L + gen}) follows.

Now suppose that $h=h_{p,r}$. We need to show that a cosingular vector $u_{rp}$ is not mapped into a~maximal submodule of $V^\HVir(h_{p,r},h_I)$. But $u_{rp}$ has $L(-p) ^{r}v$ as an additive component (see Remark~\ref{pbw}), and by construction~(\ref{Psi}), $\Psi(u_{rp})$ also must have this additive component. However, $\Psi(u_{rp})$ can not lie in a maximal $\HVir$-submodule of $V^{\HVir}(h,h_{I}) $ (see Remark~\ref{kompon}). This means that isomorphism $\Psi$ of Verma modules induces a $W(2,2) $-isomorphism of $\widetilde{L}^{W(2,2)}(h,h_{W}) $ and $L^{\HVir}(h,h_{I}) $ for all $h\in\C$. Exactness of~(\ref{seq +}) is just an application of~(\ref{seq}).
\end{proof}

\begin{Remark}
Note that the image $\Psi(u_{rp}) $ of a $W(2,2)$-cosingular vector is neither $\HVir$-singular, nor $\HVir$-cosingular in $V^{\HVir}(h_{p,r},(1+p)c_{L,I})$. For example, $L(-1) v_{0,0}$ in $V^{\HVir}(0,2c_{L,I})$ is $W(2,2)$-cosingular, but not $\HVir$-singular since $I(1) L(-1) v_{0,0}=2c_{L,I}v_{0,0}$.
\end{Remark}

If $h_{I}=(1-p) c_{L,I}$, then $\Psi $ is not an isomorphism. We shall present a $W(2,2)$-structure of Verma module later. In order to examine irreducible $W(2,2)$-modules we apply the properties of contragredient modules.

Let us recall the def\/inition of contragredient module (see~\cite{FHL}). Assume that $(M,Y_{M})$ is a~graded module over a vertex operator algebra~$V$ such that $M=\oplus _{n=0}^{\infty }M(n)$, $\dim M(n)<\infty $ and suppose that there is $\gamma \in {\C}$ such that $L(0)|M(n)\equiv (\gamma +n)\operatorname{Id}$. The contragredient module $( M^{\ast },Y_{M^{\ast}}) $ is def\/ined as follows. For every $n\in {\N}$ let $M(n)^{\ast }$ be the dual vector space and let $M^{\ast }=\oplus _{n=0}^{\infty }M(n)^{\ast }$ be a restricted dual of~$M$. Consider the natural pairing $\langle \cdot ,\cdot \rangle :M^{\ast }\otimes M\rightarrow \C$. Def\/ine the linear map $Y_{M^{\ast }}\colon V\rightarrow \operatorname{End} M^{\ast }[[z,z^{-1}]]$ such that
\begin{gather}
\langle Y_{M^{\ast }}(v,z)m^{\prime },m\rangle =\big\langle m^{\prime },Y_{M}\big(e^{zL(1)}\big({-}z^{-2}\big)^{L(0)}v,z^{-1}\big)m\big\rangle \label{contra}
\end{gather}
for each $v\in V$, $m\in M$, $m^{\prime }\in M^{\ast }$. Then $(M^{\ast },Y_{M^{\ast }}) $ is a $V$-module.

In particular, choosing $v=\omega =L_{-2}\boldsymbol{1}$ in (\ref{contra}) one gets
\begin{gather*}
\langle L(n) m^{\prime },m\rangle =\langle m^{\prime },L(-n) m\rangle .
\end{gather*}
Simple calculation with $I\in L^{\HVir}(c_{L},c_{L,I}) $ and $W\in L^{W(2,2) } ( c_{L},c_{W} ) $ shows that
\begin{gather*}
\langle I(n) m^{\prime },m\rangle = \langle m^{\prime },-I(-n) m+\delta _{n,0}2c_{L,I} \rangle , \qquad
\langle W(n) m^{\prime },m\rangle =\langle m^{\prime },W(-n) m\rangle .
\end{gather*}
Therefore we get the following result (the f\/irst and third relations were given in~\cite{nas}):

\begin{Lemma}\label{duali}
\begin{gather*}
L^{\HVir}(h,h_{I})^{\ast} \cong L^{\HVir}(h,-h_{I}+2c_{L,I}),\qquad
L^{W(2,2) }(h,h_{W})^{\ast} \cong L^{W(2,2) }(h,h_{W}) .
\end{gather*}
In particular,
\begin{gather*}
L^{\HVir}(h,(1\pm p) c_{L,I})^{\ast}\cong L^{\HVir}(h,(1\mp p) c_{L,I}).
\end{gather*}
\end{Lemma}

Directly from Theorem \ref{L za +} and Lemma \ref{duali} follows

\begin{Corollary}\label{L za -}Let $p \in \N $.
\begin{enumerate}\itemsep=0pt
\item[$(i)$] If $(h,(1-p) c_{L,I})$ is typical for $\HVir$ $($equivalently if $\big( h,\frac{1-p^{2}}{24}c_{W}\big) $ is typical for $W(2,2))$ then
\begin{gather*}
L^{\HVir} ( h,(1-p) c_{L,I} ) \cong_{W(2,2) }L^{W(2,2)}\left( h,\frac{1-p^{2}}{24}c_{W}\right) .
\end{gather*}

\item[$(ii)$] If $(h_{p,r},(1-p) c_{L,I}) \in {\mathcal {AT} }_{\HVir} ( c_L, c_{L,I} )$ $($equivalently if $ \big(h_{p,r},\frac{1-p^{2}}{24}c_{W}\big) \in {\mathcal {AT} }_{W(2,2)} ( c_L, c_W ) )$ then
\begin{gather*}
L^{\HVir} ( h_{p,r},(1-p) c_{L,I} ) \cong_{W(2,2) }\widetilde{L}^{W(2,2)}\left( h_{p,r},\frac{1-p^{2}}{24}c_{W}\right) ^{\ast}
\end{gather*}
and the short sequence of $W(2,2)$-modules
\begin{align*}
0 & \rightarrow L^{W(2,2)}\left( h_{p,r},\frac{1-p^{2}}{24}c_{W}\right)
\rightarrow L^{\HVir}\left( h_{p,r},(1-p) c_{L,I}\right) \\ %\label{seq -} \\
& \rightarrow L^{W(2,2)}\left( h_{p,r}+rp,\frac{1-p^{2}}{24}c_{W}\right)
\rightarrow0 \nonumber
\end{align*}
is exact.
\end{enumerate}
\end{Corollary}
